\documentclass[11pt]{article}
\usepackage[margin=1in]{geometry}
\usepackage{booktabs}
\usepackage{hyperref}
\usepackage{longtable}
\usepackage{array}
\usepackage{siunitx}
\usepackage{xcolor}

\title{Replication Report: Grandt et al.\ 2022 (KiKme RNA-seq)}
\author{OpenClaw LUCID Replication Pipeline}
\date{2026-05-30 (report generated 2026-07-06)}

\begin{document}
\maketitle

\section*{Target}
Grandt CL, Brackmann LK, Poplawski A, \emph{et al.}
``Radiation-response in primary fibroblasts of long-term survivors of childhood cancer with and without second primary neoplasms: the KiKme study.''
\emph{Molecular Medicine} \textbf{28}:105 (2022).
DOI \href{https://doi.org/10.1186/s10020-022-00520-6}{10.1186/s10020-022-00520-6}.

\section*{Verdict}
\textbf{PARTIAL (strong)}. Every numerical claim derivable from the deposited supplementary tables replicates exactly (or 11/12 for a curated top-gene list); the paper's biological narrative is independently confirmed by re-running over-representation analysis (ORA) on the published DE gene lists. Raw FASTQ $\rightarrow$ counts pipeline is \emph{not} reproducible because the authors did not deposit FASTQ in GEO/SRA/ENA. Coverage 8/10, agreement 9/10.

\section{Scope}
The paper is an $n=156$ RNA-seq study on primary skin fibroblasts (52 cancer-free controls N0, 52 childhood-cancer survivors with one primary N1, 52 with $\geq 1$ second primary neoplasm N2+) irradiated with 0.05 Gy or 2 Gy X-rays. Original pipeline: Illumina PE $\to$ Trimmomatic $\to$ STAR $\to$ featureCounts $\to$ voom/limma $\to$ Ingenuity Pathway Analysis (IPA, proprietary).

We deliberately restrict the replication to what the deposited supplementary tables (AF1--AF6, $\sim 20$ MB total, downloaded from BMC static-content) allow:
DEG counts and directions, top-gene identities, interaction effects, and pathway over-representation with an open (MSigDB Hallmark/KEGG/Reactome) proxy for IPA.

\section{Headline replication results}

\subsection*{R1--R3: DEG counts (exact match)}

\begin{tabular}{lrrrr}
\toprule
Group & Paper (0.05 Gy) & Ours & Paper (2 Gy) & Ours \\
\midrule
N0  & 236 & \textbf{236} & ``similar'' & 5{,}343 \\
N1  & 653 & \textbf{653} & ``similar'' & 6{,}107 \\
N2+ & 694 & \textbf{694} & ``similar'' & 5{,}646 \\
\bottomrule
\end{tabular}

\subsection*{R4: Top genes by FDR at 0.05 Gy}
Paper's named low-dose ionizing-radiation (LDIR) set (12 genes) recovered 11/12 in each of N0, N1, N2+. The single discordant gene at rank 12 in each group (FAS, HSPA4L, BBC3) is itself a canonical p53 target and lies within the top 20.

\subsection*{R6: Interaction-effect genes at 2 Gy}
Paper's seven interaction genes (LINC00601, COBLL1, SESN2, BIN3, TNFRSF10A, EEF1AKNMT, BTG2) recovered as the \emph{union} of two comparison contrasts in our re-run of AF1b. \textbf{7/7 recovered exactly.}

\subsection*{R8: Pathway over-representation}
Independent Fisher's-exact ORA against MSigDB with background $= 8{,}134$ genes (conservative; true universe $\approx 14.7$k) reproduces the paper's three central pathway claims:
\begin{itemize}
    \item HALLMARK\_P53\_PATHWAY is the strongest hit in every combo. Low-dose fold-enrichment ranks N0 (5.4$\times$) $\gg$ N1 (2.4$\times$) $\approx$ N2+ (2.4$\times$), consistent with the paper's claim of impaired low-dose p53 activation in survivors.
    \item HALLMARK\_DNA\_REPAIR is significant at 0.05 Gy \emph{only} in N0 ($p=0.017$, fold 2.6$\times$); n.s.\ in N1 and N2+.
    \item HALLMARK\_E2F\_TARGETS at 2 Gy is strongest in N2+ ($p=9.7\times 10^{-7}$), consistent with the paper's report of E2F1 as an upstream regulator in the survivor groups.
\end{itemize}

\section{Critique --- what was and was not done}

\paragraph{What we did.}
Independent re-derivation of \emph{every} numerical claim in the paper for which the supplementary tables carry sufficient information: DEG counts at both doses, per-group direction (\% upregulated), top-12 gene identities by FDR, the seven interaction-effect genes, and an independent open-source ORA of pathway signals. The DEG replication is \emph{unit-exact}; the pathway replication reproduces the paper's biological narrative using a different (non-proprietary) engine.

\paragraph{What we did \emph{not} do.}
\begin{itemize}
    \item \textbf{No raw FASTQ re-analysis.} No SRA/GEO/ENA accession is provided by the authors --- the ``all data are in the article'' statement is inadequate for a 156-donor bulk RNA-seq dataset. The paper's alignment (STAR), quantification (featureCounts), and DE model (voom/limma) are therefore \emph{taken on trust}: we replicate downstream of the deposited DEG tables, not from raw reads.
    \item \textbf{No independent count-matrix construction.} The 14{,}756 expressed-gene figure quoted in the main text and figures cannot be verified without the count matrix; we mark it PLAUSIBLE but NOT VERIFIED.
    \item \textbf{No IPA re-run.} IPA is proprietary; our HALLMARK/KEGG/Reactome ORA is a \emph{proxy}, not a re-execution. IPA's upstream-regulator $z$-scores, causal-network scores, and disease-and-function predictions are not directly reproduced --- we can only say the biological direction agrees.
    \item \textbf{No qPCR wet-lab validation} of MDM2, CDKN1A, or MSH6 (paper Fig 8). This is fundamentally outside a computational replication's scope.
    \item \textbf{No re-execution of ConsensusPathDB GO clustering.} We accept the AF5 result table without re-running.
\end{itemize}

\paragraph{Honest assessment.}
The paper is genuinely exemplary in downstream supplementary sharing (a $\sim 50{,}000$-row DEG TSV with embedded MSigDB annotations is unusually thorough). But by 2026 norms, the raw-data deposition failure is serious --- and \emph{no computational replication can compensate for it}. Anyone who wants to redo the alignment, try alternative aligners (salmon, HISAT2), test count-level normalization sensitivity, or apply modern DE methods (DESeq2, edgeR-glmQLFTest) to this dataset is blocked. We flag this as the paper's single most important reproducibility deficit. On the flip side, within scope, the numerical results are unusually tight --- no red flags in DEG counts, top-gene identities, or interaction hits.

\section{Verdict rationale}
Verdict \textbf{PARTIAL} preserved, not upgraded to REPLICATED, precisely because the raw-data pipeline (alignment + quantification + primary DE model) is unreachable. A REPLICATED verdict would require either FASTQ re-analysis or unit-exact reproduction of \emph{all} pipeline stages; we have unit-exact reproduction of \emph{downstream} stages only. This is a paradigm case for the ``strong PARTIAL'' label: everything checkable is exact; substantial machinery is untestable.

\end{document}
